\setgebiet{6}{Numerik}{tief}
\chapter{Numerik}
\marginnote{\small\itshape\color{grau}\textbf{Sie brauchen:} Band 5, Kapitel 1--14. Aus Band 2 Matrizen, LR, QR, Eigenwerte, SVD (Gebiet 2); aus Band 3 Taylor, Mittelwertsatz, Integrale (Gebiet 3). Ein Taschenrechner oder ein paar Zeilen Python helfen bei den längeren Rechnungen.}[12pt]
\worum{Numerik ist Mathematik unter der Bedingung, dass man nur endlich viele Stellen hat und nur endlich viele Rechenschritte. Die Aufgaben üben, was dadurch neu wird: Fehler erkennen, bevor sie entstehen (Kondition, Stabilität), Verfahren vergleichen (Aufwand, Konvergenzordnung), und Ergebnisse prüfen, die man nicht exakt kennt.}

\section{Was Sie hier üben}
\begin{center}
{\footnotesize\sffamily\setlength{\tabcolsep}{4pt}
\renewcommand{\arraystretch}{1.5}
\begin{tabular}{@{}lP{4.9cm}lcll@{}}
\toprule
\textbf{Nr.} & \textbf{Titel} & \textbf{Sorte} & \textbf{Grad} & \textbf{Zeit} & \textbf{Band/Kap.}\\
\midrule
5.1 & Binär, Epsilon, $0{,}1 + 0{,}2$ & \typtypisch & \grad1 & 10 min & 5/1\\
5.2 & Gleich ist nicht gleich & \typstolper & \grad1 & 10 min & 5/1\\
5.3 & $-6{,}25$ als double & \typwerk & \grad1 & 10 min & 5/1\\
5.4 & Die Mitternachtsformel versagt & \typwerk & \grad2 & 15 min & 5/1\\
5.5 & $\frac{1 - \cos x}{x^{2}}$ für kleines $x$ & \typwerk & \grad2 & 10 min & 5/1\\
5.6 & Zehn Millionen mal $0{,}1$ & \typwerk & \grad2 & 10 min & 5/1\\
5.7 & Die Kondition elementarer Funktionen & \typtypisch & \grad2 & 15 min & 5/2\\
5.8 & Kondition einer doppelten Nullstelle & \typwerk & \grad2 & 15 min & 5/2\\
5.9 & Eine Konditionszahl von $40\,000$ & \typwerk & \grad2 & 15 min & 5/2\\
5.10 & Das Skalarprodukt ist rückwärtsstabil & \typwerk & \grad3 & 20 min & 5/2\\
5.11 & Was kostet es? & \typtypisch & \grad2 & 15 min & 5/3\\
5.12 & Eine LR-Zerlegung mit Probe & \typwerk & \grad2 & 15 min & 5/3\\
5.13 & Cholesky mit großen Zahlen & \typwerk & \grad1 & 10 min & 5/3\\
5.14 & QR mit Householder & \typwerk & \grad2 & 20 min & 5/4\\
5.15 & Eine Ausgleichsgerade & \typwerk & \grad1 & 10 min & 5/4\\
5.16 & Jacobi divergiert & \typstolper & \grad2 & 10 min & 5/5\\
5.17 & Jacobi und Gauß--Seidel & \typwerk & \grad2 & 20 min & 5/5\\
5.18 & Zwei CG-Schritte & \typwerk & \grad3 & 20 min & 5/5\\
5.19 & Gerschgorin und Rayleigh & \typtypisch & \grad2 & 10 min & 5/6\\
5.20 & Die Potenzmethode & \typwerk & \grad2 & 15 min & 5/6\\
5.21 & Eine SVD von Hand & \typwerk & \grad2 & 15 min & 5/6\\
5.22 & Bisektion gegen Newton & \typtypisch & \grad2 & 15 min & 5/7\\
5.23 & Newton im Kreis & \typstolper & \grad2 & 15 min & 5/7\\
5.24 & $\sqrt2$ auf zehn Stellen & \typwerk & \grad1 & 10 min & 5/7\\
5.25 & $x = \cos x$: Fixpunkt oder Newton & \typwerk & \grad2 & 15 min & 5/7\\
5.26 & Newton in zwei Dimensionen & \typwerk & \grad3 & 20 min & 5/7\\
5.27 & Der Interpolationsfehler & \typtypisch & \grad2 & 15 min & 5/8\\
5.28 & Höherer Grad, schlechteres Ergebnis & \typstolper & \grad2 & 15 min & 5/8\\
5.29 & Das Newton-Schema & \typwerk & \grad1 & 10 min & 5/8\\
5.30 & Tschebyschow-Knoten & \typwerk & \grad2 & 15 min & 5/8\\
5.31 & Ein natürlicher Spline & \typwerk & \grad2 & 20 min & 5/9\\
5.32 & Eine DFT von Hand & \typwerk & \grad2 & 10 min & 5/9\\
5.33 & Beste Quadratmittel-Approximation & \typwerk & \grad2 & 15 min & 5/10\\
5.34 & Beste gleichmäßige Approximation & \typwerk & \grad2 & 15 min & 5/10\\
5.35 & Die Fehlerordnung messen & \typtypisch & \grad2 & 15 min & 5/11\\
5.36 & Trapez und Simpson & \typwerk & \grad2 & 15 min & 5/11\\
5.37 & Das Romberg-Schema & \typwerk & \grad2 & 15 min & 5/11\\
5.38 & Gauß--Legendre mit drei Punkten & \typwerk & \grad2 & 15 min & 5/11\\
5.39 & Euler explodiert & \typstolper & \grad2 & 15 min & 5/12\\
5.40 & Euler und Runge--Kutta & \typwerk & \grad2 & 20 min & 5/12\\
5.41 & Ein Randwertproblem auf drei Punkten & \typwerk & \grad2 & 15 min & 5/13\\
5.42 & Gradientenverfahren im Tal & \typtypisch & \grad2 & 15 min & 5/14\\
5.43 & Ein Newton-Schritt auf Rosenbrock & \typwerk & \grad2 & 15 min & 5/14\\
5.44 & $\pi$ mit Monte-Carlo & \typwerk & \grad2 & 10 min & 5/14\\
\bottomrule
\end{tabular}}
\end{center}

\section{Woran Sie merken, dass Sie es können}
\begin{itemize}[leftmargin=1.2em,itemsep=2pt]
\item Sie sehen einer Formel an, wo Auslöschung droht, und kennen die Umformung, die sie vermeidet.
\item Sie unterscheiden Kondition (Eigenschaft des Problems) von Stabilität (Eigenschaft des Verfahrens) und wissen, welche von beiden man verbessern kann.
\item Sie können die Konvergenzordnung eines Verfahrens aus zwei Fehlern ablesen --- und wissen, dass Newton quadratisch konvergiert, wenn er konvergiert.
\item Sie kennen zu jedem Standardverfahren ein Beispiel, an dem es scheitert: Jacobi ohne Diagonaldominanz, Newton ohne guten Start, Interpolation auf gleichabständigen Knoten, expliziter Euler bei steifen Gleichungen.
\end{itemize}

\section{Aufgaben}

%%%%%%%%%%% Kapitel 1
\stupser{tief}{$0{,}1 = \frac1{10}$, und $10 = 2\cdot5$. Ein Bruch ist im Binärsystem genau dann endlich, wenn der Nenner eine Zweierpotenz ist.}{22pt}
\begin{karte}{tief}{5.1}{\typtypisch}{\meta{1}{10 min}{Band 5, Kap.\ 1}}
\kartentitel{Binär, Epsilon, $0{,}1 + 0{,}2$}
\worum{Drei Tatsachen, die jeder Programmierer einmal schmerzhaft lernt --- hier mit Begründung.}
\methoden{Gleitpunktzahlen, Binärdarstellung, Maschinenepsilon}
\begin{enumerate}[label=(\alph*),leftmargin=1.6em,itemsep=2pt]
\item Schreiben Sie $0{,}1$ im Binärsystem. Warum ist die Darstellung nicht endlich?
\item Erklären Sie, warum ein Rechner $0{,}1 + 0{,}2 = 0{,}30000000000000004$ ausgibt.
\item Das Maschinenepsilon $\varepsilon_M$ von double ist $2^{-52}\approx2{,}2\cdot10^{-16}$. Was bedeutet es? Was liefert der Rechner für $1 + 10^{-16}$, was für $1 + 10^{-15}$?
\item Wie viele Dezimalstellen sind in double verlässlich, wie viele in single ($\varepsilon_M = 2^{-23}$)?
\end{enumerate}
\end{karte}

\stupser{akzent}{$10^{16} + 1$ liegt zwischen zwei benachbarten double-Zahlen; der Abstand dort ist $2$. Also $10^{16} + 1 = 10^{16}$ im Rechner.}{22pt}
\begin{karte}{akzent}{5.2}{\typstolper}{\meta{1}{10 min}{Band 5, Kap.\ 1}}
\kartentitel{Gleich ist nicht gleich}
\worum{Gleitpunktzahlen mit \texttt{==} vergleichen, auf das Assoziativgesetz vertrauen, Summen in beliebiger Reihenfolge bilden --- drei Dinge, die bei reellen Zahlen richtig und im Rechner falsch sind.}
\methoden{Gleitpunktzahlen, Absorption, Summationsreihenfolge}
\begin{enumerate}[label=(\alph*),leftmargin=1.6em,itemsep=2pt]
\item Was liefert der Rechner für $(10^{16} + 1) - 10^{16}$? Und für $10^{16} + (1 - 10^{16})$?
\item Was liefert $(0{,}1 + 0{,}2) + 0{,}3$, was $0{,}1 + (0{,}2 + 0{,}3)$? Welches Rechengesetz ist verletzt?
\item Eine Schleife soll bei $x = 1$ abbrechen und zählt $x$ in Schritten von $0{,}1$ hoch. Warum endet sie nie? Wie prüft man Gleichheit richtig?
\item Man summiert $1 + 10^{-16} + 10^{-16} + \dots$ ($10^{4}$ Summanden $10^{-16}$). Vorwärts summiert kommt $1$ heraus. Was kommt heraus, wenn man die kleinen Zahlen zuerst addiert? Welche Reihenfolge ist allgemein günstiger?
\end{enumerate}
\end{karte}

\stupser{gruen}{$6{,}25 = (110{,}01)_2 = (1{,}1001)_2\cdot2^{2}$. Exponent $2$ plus Bias $1023$ ist $1025$.}{22pt}
\begin{karte}{gruen}{5.3}{\typwerk}{\meta{1}{10 min}{Band 5, Kap.\ 1, Aufgabe 1.1}}
\kartentitel{$-6{,}25$ als double}
\worum{Vorzeichen, Exponent, Mantisse: einmal von Hand, damit man weiß, was in den 64 Bit steht.}
\methoden{Binärdarstellung, Gleitpunktzahlen}
Stellen Sie $x = -6{,}25$ im Format double dar (Vorzeichenbit, 11 Exponentenbits mit Bias $1023$, 52 Mantissenbits).
\end{karte}

\stupser{gruen}{$x_2 = \frac{p - \sqrt{p^{2} - 4q}}2$ subtrahiert zwei fast gleiche Zahlen. Vieta: $x_1x_2 = q$.}{22pt}
\begin{karte}{gruen}{5.4}{\typwerk}{\meta{2}{15 min}{Band 5, Kap.\ 1, Aufgabe 1.2}}
\kartentitel{Die Mitternachtsformel versagt}
\worum{Das Lehrbuchbeispiel für Auslöschung: eine richtige Formel, falsch ausgewertet. Und die Umformung, die alles rettet.}
\methoden{Auslöschung, Vieta, Umformen vor dem Rechnen}
Berechnen Sie die Nullstellen von $x^{2} - 10^{8}x + 1 = 0$ mit der Mitternachtsformel in doppelter Genauigkeit. Was geht schief, und wie behebt man es?
\end{karte}

\stupser{gruen}{$\cos(10^{-8}) = 1 - 5\cdot10^{-17}$ wird zu $1$ gerundet. Umformung: $1 - \cos x = 2\sin^{2}\frac x2$.}{22pt}
\begin{karte}{gruen}{5.5}{\typwerk}{\meta{2}{10 min}{Band 5, Kap.\ 1, Aufgabe 1.3}}
\kartentitel{$\frac{1 - \cos x}{x^{2}}$ für kleines $x$}
\worum{Noch einmal Auslöschung, diesmal mit einer trigonometrischen Identität als Rettung --- oder mit der Taylorreihe.}
\methoden{Auslöschung, Trigonometrische Identität, Taylorreihe}
Warum liefert $\frac{1 - \cos x}{x^{2}}$ für $x = 10^{-8}$ den Wert $0$, und wie rechnet man richtig?
\end{karte}

\stupser{gruen}{Jede Addition kostet einen relativen Fehler $\le\varepsilon_M$; nach $n$ Additionen ist der Gesamtfehler bis zu $n\varepsilon_M$ mal die Summe der Beträge.}{22pt}
\begin{karte}{gruen}{5.6}{\typwerk}{\meta{2}{10 min}{Band 5, Kap.\ 1, Aufgabe 1.4}}
\kartentitel{Zehn Millionen mal $0{,}1$}
\worum{Rundungsfehler sammeln sich an. Wie schnell, hängt vom Verfahren ab --- naive Summation ist das schlechteste.}
\methoden{Rundungsfehler, Fehlerschranke, Summationsreihenfolge}
Man summiere $10^{7}$-mal die Zahl $0{,}1$ in einfacher Genauigkeit ($\varepsilon_M\approx6\cdot10^{-8}$). Welchen Fehler erwartet man, und welchen bekommt man tatsächlich? Wie macht man es besser?
\end{karte}

%%%%%%%%%%% Kapitel 2
\stupser{tief}{Relative Kondition $\kappa(x) = \bigl|\frac{x\,f'(x)}{f(x)}\bigr|$: Um welchen Faktor wird ein relativer Eingabefehler verstärkt?}{22pt}
\begin{karte}{tief}{5.7}{\typtypisch}{\meta{2}{15 min}{Band 5, Kap.\ 2}}
\kartentitel{Die Kondition elementarer Funktionen}
\worum{Kondition ist eine Eigenschaft des Problems, nicht des Rechners. Hier für die Funktionen, aus denen alle anderen gebaut sind.}
\methoden{Konditionszahl, Relative Fehler}
Bestimmen Sie die relative Konditionszahl und sagen Sie, wo sie groß wird:
\begin{enumerate}[label=(\alph*),leftmargin=1.6em,itemsep=2pt]
\item $f(x) = \sqrt x$, \quad $f(x) = e^{x}$, \quad $f(x) = \ln x$,
\item $f(x) = \sin x$,
\item $f(x, y) = x - y$ bezüglich $x$ (bei festem $y$),
\item Welche der Funktionen ist überall gut konditioniert? Warum ist die Subtraktion die gefährlichste Grundrechenart?
\end{enumerate}
\end{karte}

\stupser{gruen}{Nullstellen $1\pm\delta$. Stört man den konstanten Koeffizienten $c = 1 - \delta^{2}$, so ist $\frac{dx_1}{dc} = -\frac1{2\delta}$.}{22pt}
\begin{karte}{gruen}{5.8}{\typwerk}{\meta{2}{15 min}{Band 5, Kap.\ 2, Aufgabe 2.1}}
\kartentitel{Kondition einer doppelten Nullstelle}
\worum{Nahe beieinanderliegende Nullstellen sind schlecht konditioniert --- und bei $\delta\to0$ wird die Kondition unendlich. Hier mit Formel.}
\methoden{Konditionszahl, Ableiten nach dem Parameter}
Bestimmen Sie die relative Kondition der Nullstellen von $p(x) = x^{2} - 2x + (1 - \delta^{2})$ in Abhängigkeit von $\delta$ und deuten Sie das Ergebnis.
\end{karte}

\stupser{gruen}{$\operatorname{cond}_\infty(A) = \|A\|_\infty\|A^{-1}\|_\infty$; $A^{-1} = \frac1{0{,}0001}\begin{pmatrix}1{,}0001 & -1\\ -1 & 1\end{pmatrix}$.}{22pt}
\begin{karte}{gruen}{5.9}{\typwerk}{\meta{2}{15 min}{Band 5, Kap.\ 2, Aufgabe 2.2}}
\kartentitel{Eine Konditionszahl von $40\,000$}
\worum{Eine Störung von $10^{-4}$ in der rechten Seite ändert die Lösung um $100\,\%$. Die Konditionszahl sagt das vorher.}
\methoden{Konditionszahl, Störungsabschätzung, Matrixnorm}
$A = \begin{pmatrix}1 & 1\\ 1 & 1{,}0001\end{pmatrix}$, $b = \begin{pmatrix}2\\ 2{,}0001\end{pmatrix}$.
\begin{enumerate}[label=(\alph*),leftmargin=1.6em,itemsep=2pt]
\item Lösen Sie $Ax = b$ exakt.
\item Stören Sie $b$ zu $\tilde b = (2;\ 2{,}0002)^{\mathsf T}$ und lösen Sie erneut.
\item Berechnen Sie $\operatorname{cond}_\infty(A)$ und prüfen Sie die Störungsabschätzung $\frac{\|\Delta x\|}{\|x\|}\le\operatorname{cond}(A)\frac{\|\Delta b\|}{\|b\|}$.
\end{enumerate}
\end{karte}

\stupser{gruen}{Jeder Term $x_iy_i$ wird von höchstens $n$ Rundungen getroffen: $\tilde s = \sum x_iy_i(1 + \theta_i)$ mit $|\theta_i|\le n\varepsilon_M$. Das ist das exakte Skalarprodukt leicht gestörter $x_i$.}{22pt}
\begin{karte}{gruen}{5.10}{\typwerk}{\meta{3}{20 min}{Band 5, Kap.\ 2, Aufgabe 2.3}}
\kartentitel{Das Skalarprodukt ist rückwärtsstabil}
\worum{Rückwärtsanalyse: Das berechnete Ergebnis ist das exakte Ergebnis für leicht veränderte Eingaben. Hier der einfachste Fall, Schritt für Schritt.}
\methoden{Rückwärtsanalyse, Standardmodell der Rundung}
Zeigen Sie, dass die sequenzielle Berechnung von $s = \sum_{i=1}^{n}x_iy_i$ rückwärtsstabil ist: $\tilde s = \sum\tilde x_iy_i$ mit $|\tilde x_i - x_i|\le n\varepsilon_M|x_i|$ (bis auf $O(\varepsilon_M^{2})$).
\end{karte}

%%%%%%%%%%% Kapitel 3
\stupser{tief}{LR-Zerlegung $\approx\frac{2}{3}n^{3}$ Gleitpunktoperationen, Vorwärts- und Rückwärtseinsetzen je $\approx n^{2}$, Inverse $\approx2n^{3}$, Cramer $\sim(n + 1)!$.}{22pt}
\begin{karte}{tief}{5.11}{\typtypisch}{\meta{2}{15 min}{Band 5, Kap.\ 3}}
\kartentitel{Was kostet es?}
\worum{Operationen zählen ist die erste Pflicht der Numerik. Die Antworten erklären, warum man nie die Inverse berechnet und nie Cramer benutzt.}
\methoden{Aufwand zählen, Größenordnung schätzen}
\begin{enumerate}[label=(\alph*),leftmargin=1.6em,itemsep=2pt]
\item Zählen Sie die Multiplikationen in der LR-Zerlegung einer $n\times n$-Matrix (führender Term) und im Rückwärtseinsetzen.
\item Ein Rechner schafft $10^{9}$ Operationen pro Sekunde. Wie lange dauert $Ax = b$ für $n = 1000$, $10^{4}$, $10^{5}$?
\item Man will $Ax = b$ für $100$ rechte Seiten lösen. Vergleichen Sie: (i) $100$-mal LR, (ii) einmal LR, $100$-mal einsetzen, (iii) $A^{-1}$ berechnen und $100$-mal multiplizieren.
\item Die Cramersche Regel mit Determinanten per Entwicklung braucht etwa $(n + 1)!$ Operationen. Für welches $n$ überschreitet das das Alter des Universums ($4\cdot10^{17}$ s)?
\end{enumerate}
\end{karte}

\stupser{gruen}{Pivot $2$, Multiplikatoren $\ell_{21} = -1{,}5$, $\ell_{31} = -1$; danach $\ell_{32} = 4$. Die Multiplikatoren \emph{sind} $L$.}{22pt}
\begin{karte}{gruen}{5.12}{\typwerk}{\meta{2}{15 min}{Band 5, Kap.\ 3, Aufgabe 3.1}}
\kartentitel{Eine LR-Zerlegung mit Probe}
\worum{Gauß als Zerlegung, dann zwei Dreieckssysteme, dann die Probe $LR = A$.}
\methoden{LR-Zerlegung, Rückwärtseinsetzen, Probe}
Berechnen Sie die LR-Zerlegung von $A = \begin{pmatrix}2 & 1 & -1\\ -3 & -1 & 2\\ -2 & 1 & 2\end{pmatrix}$ (ohne Pivotsuche) und lösen Sie damit $Ax = (8, -11, -3)^{\mathsf T}$.
\end{karte}

\stupser{gruen}{$\ell_{11} = 2$, $\ell_{21} = 6$, $\ell_{31} = -8$; $\ell_{22} = \sqrt{37 - 36} = 1$.}{22pt}
\begin{karte}{gruen}{5.13}{\typwerk}{\meta{1}{10 min}{Band 5, Kap.\ 3, Aufgabe 3.2}}
\kartentitel{Cholesky mit großen Zahlen}
\worum{Die Zerlegung geht glatt auf --- ein Zeichen, dass die Aufgabe konstruiert ist, und eine gute Übung für die Spaltenformel.}
\methoden{Cholesky-Zerlegung, Positive Definitheit}
Bestimmen Sie die Cholesky-Zerlegung von $A = \begin{pmatrix}4 & 12 & -16\\ 12 & 37 & -43\\ -16 & -43 & 98\end{pmatrix}$ und prüfen Sie $LL^{\mathsf T} = A$.
\end{karte}

%%%%%%%%%%% Kapitel 4
\stupser{gruen}{Erste Spalte $(3, 4, 0)$, Norm $5$, $v = (3 - 5, 4, 0) = (-2, 4, 0)$. $H_1 = I - 2\frac{vv^{\mathsf T}}{v^{\mathsf T}v}$ wirkt nur auf die ersten beiden Komponenten.}{22pt}
\begin{karte}{gruen}{5.14}{\typwerk}{\meta{2}{20 min}{Band 5, Kap.\ 4, Aufgabe 4.1}}
\kartentitel{QR mit Householder}
\worum{Zwei Spiegelungen machen aus einer $3\times2$-Matrix eine obere Dreiecksmatrix. Die Zahlen sind so gewählt, dass die erste Spiegelung die aus Gebiet 2 ist.}
\methoden{Householder-Spiegelung, QR-Zerlegung}
Bestimmen Sie die QR-Zerlegung von $A = \begin{pmatrix}3 & 1\\ 4 & 2\\ 0 & 5\end{pmatrix}$ mit Householder-Spiegelungen.
\end{karte}

\stupser{gruen}{Normalengleichungen $A^{\mathsf T}A\begin{pmatrix}\alpha\\ \beta\end{pmatrix} = A^{\mathsf T}b$ mit $A^{\mathsf T}A = \begin{pmatrix}4 & 10\\ 10 & 30\end{pmatrix}$, $A^{\mathsf T}b = \begin{pmatrix}20\\ 59{,}7\end{pmatrix}$.}{22pt}
\begin{karte}{gruen}{5.15}{\typwerk}{\meta{1}{10 min}{Band 5, Kap.\ 4, Aufgabe 4.2}}
\kartentitel{Eine Ausgleichsgerade}
\worum{Kleinste Quadrate mit vier Punkten; dieselbe Rechnung wie Aufgabe 2.41, hier mit der Frage, wann man die Normalengleichungen \emph{nicht} benutzen sollte.}
\methoden{Kleinste Quadrate, Normalengleichungen, Konditionszahl}
Passen Sie $y = \alpha + \beta t$ an $(1;\ 2{,}1)$, $(2;\ 3{,}9)$, $(3;\ 6{,}2)$, $(4;\ 7{,}8)$ an. Berechnen Sie das Residuum. Warum ist für große, schlecht konditionierte Probleme die QR-Zerlegung den Normalengleichungen vorzuziehen?
\end{karte}

%%%%%%%%%%% Kapitel 5
\stupser{akzent}{Die Iterationsmatrix von Jacobi ist $-D^{-1}(L + U)$; für $\begin{pmatrix}1 & 2\\ 3 & 1\end{pmatrix}$ ist das $\begin{pmatrix}0 & -2\\ -3 & 0\end{pmatrix}$ mit Eigenwerten $\pm\sqrt6$.}{22pt}
\begin{karte}{akzent}{5.16}{\typstolper}{\meta{2}{10 min}{Band 5, Kap.\ 5}}
\kartentitel{Jacobi divergiert}
\worum{Iterative Verfahren konvergieren nicht immer. Diagonaldominanz ist die Bedingung, die man prüft, bevor man loslegt --- und sie hängt von der Zeilenreihenfolge ab.}
\methoden{Spektralradius, Diagonaldominanz, Gegenbeispiel}
$A = \begin{pmatrix}1 & 2\\ 3 & 1\end{pmatrix}$, $b = \begin{pmatrix}3\\ 4\end{pmatrix}$ (Lösung $(1, 1)$).
\begin{enumerate}[label=(\alph*),leftmargin=1.6em,itemsep=2pt]
\item Führen Sie drei Jacobi-Schritte ab $x^{(0)} = 0$ aus. Was passiert?
\item Bestimmen Sie den Spektralradius der Iterationsmatrix und erklären Sie damit (a).
\item Vertauschen Sie die beiden Gleichungen und wiederholen Sie (a), (b).
\item Warum konvergiert das Verfahren bei strikter Diagonaldominanz immer? (Skizze mit der Zeilensummennorm.)
\end{enumerate}
\end{karte}

\stupser{gruen}{Jacobi: $x_1^{(k+1)} = \frac{2 + x_2^{(k)}}4$, $x_2^{(k+1)} = \frac{4 + x_1^{(k)} + x_3^{(k)}}4$, $x_3^{(k+1)} = \frac{10 + x_2^{(k)}}4$. Gauß--Seidel benutzt die neuen Werte sofort.}{22pt}
\begin{karte}{gruen}{5.17}{\typwerk}{\meta{2}{20 min}{Band 5, Kap.\ 5, Aufgabe 5.1}}
\kartentitel{Jacobi und Gauß--Seidel}
\worum{Zwei Iterationen auf derselben tridiagonalen Matrix. Gauß--Seidel ist etwa doppelt so schnell --- und man sieht, warum.}
\methoden{Jacobi-Verfahren, Gauß–Seidel-Verfahren, Konvergenzrate}
Lösen Sie $\begin{pmatrix}4 & -1 & 0\\ -1 & 4 & -1\\ 0 & -1 & 4\end{pmatrix}x = \begin{pmatrix}2\\ 4\\ 10\end{pmatrix}$ mit Jacobi und Gauß--Seidel, Startwert $0$, je vier Schritte. Vergleichen Sie die Fehler.
\end{karte}

\stupser{gruen}{$r_0 = b - Ax_0 = (-8, -3)$, $p_0 = r_0$, $\alpha_0 = \frac{r_0^{\mathsf T}r_0}{p_0^{\mathsf T}Ap_0}$. Nach zwei Schritten muss $r_2 = 0$ sein.}{22pt}
\begin{karte}{gruen}{5.18}{\typwerk}{\meta{3}{20 min}{Band 5, Kap.\ 5, Aufgabe 5.2}}
\kartentitel{Zwei CG-Schritte}
\worum{Das Verfahren der konjugierten Gradienten löst ein $n\times n$-System in höchstens $n$ Schritten exakt. Hier $n = 2$, jede Zahl von Hand.}
\methoden{Konjugierte Gradienten, Krylow-Raum}
Führen Sie zwei CG-Schritte für $A = \begin{pmatrix}4 & 1\\ 1 & 3\end{pmatrix}$, $b = \begin{pmatrix}1\\ 2\end{pmatrix}$, $x_0 = \begin{pmatrix}2\\ 1\end{pmatrix}$ aus und prüfen Sie, dass die exakte Lösung erreicht ist.
\end{karte}

%%%%%%%%%%% Kapitel 6
\stupser{tief}{Gerschgorin: Jeder Eigenwert liegt in einem Kreis um ein Diagonalelement mit Radius gleich der Summe der übrigen Beträge der Zeile.}{22pt}
\begin{karte}{tief}{5.19}{\typtypisch}{\meta{2}{10 min}{Band 5, Kap.\ 6}}
\kartentitel{Gerschgorin und Rayleigh}
\worum{Zwei Werkzeuge, mit denen man Eigenwerte eingrenzt, ohne sie zu berechnen --- und mit denen man prüft, ob ein berechneter Eigenwert plausibel ist.}
\methoden{Eigenwerte einschließen}
$A = \begin{pmatrix}2 & 1\\ 1 & 3\end{pmatrix}$ (Eigenwerte $\frac{5\pm\sqrt5}2$).
\begin{enumerate}[label=(\alph*),leftmargin=1.6em,itemsep=2pt]
\item Geben Sie die Gerschgorin-Kreise an. Welches Intervall liefern sie für die Eigenwerte?
\item Berechnen Sie den Rayleigh-Quotienten $\frac{x^{\mathsf T}Ax}{x^{\mathsf T}x}$ für $x = (1, 1)$, $x = (1, 2)$ und $x = (1;\ 1{,}618)$. Was beobachten Sie?
\item Warum liegt der Rayleigh-Quotient bei symmetrischen Matrizen immer zwischen dem kleinsten und dem größten Eigenwert?
\item Wie gut ist der Rayleigh-Quotient, wenn $x$ um $\varepsilon$ vom Eigenvektor abweicht?
\end{enumerate}
\end{karte}

\stupser{gruen}{$x_{k+1} = Ax_k/\|Ax_k\|_\infty$; der Rayleigh-Quotient oder der Quotient der größten Komponenten nähert $\lambda_1 = 3{,}618$.}{22pt}
\begin{karte}{gruen}{5.20}{\typwerk}{\meta{2}{15 min}{Band 5, Kap.\ 6, Aufgabe 6.1}}
\kartentitel{Die Potenzmethode}
\worum{Immer wieder mit $A$ multiplizieren: Der größte Eigenwert setzt sich durch. Wie schnell, bestimmt das Verhältnis $|\lambda_2/\lambda_1|$.}
\methoden{Potenzmethode, Eigenwerte einschließen, Konvergenzrate}
Bestimmen Sie den betragsgrößten Eigenwert von $A = \begin{pmatrix}2 & 1\\ 1 & 3\end{pmatrix}$ mit der Potenzmethode, Startvektor $(1, 0)^{\mathsf T}$, sechs Schritte. Wie schnell konvergiert das Verfahren, und warum?
\end{karte}

\stupser{gruen}{$A^{\mathsf T}A = \begin{pmatrix}25 & 20\\ 20 & 25\end{pmatrix}$: Eigenwerte $45$ und $5$. Dann $u_i = Av_i/\sigma_i$.}{22pt}
\begin{karte}{gruen}{5.21}{\typwerk}{\meta{2}{15 min}{Band 5, Kap.\ 6, Aufgabe 6.2}}
\kartentitel{Eine SVD von Hand}
\worum{Dieselbe Matrix wie in Aufgabe 2.44 --- hier mit der numerischen Frage: Warum rechnet man die SVD in der Praxis \emph{nicht} über $A^{\mathsf T}A$?}
\methoden{Singulärwertzerlegung, Konditionszahl}
Berechnen Sie die SVD von $A = \begin{pmatrix}3 & 0\\ 4 & 5\end{pmatrix}$, geben Sie $\operatorname{cond}_2(A)$ an, und erklären Sie, warum der Umweg über $A^{\mathsf T}A$ numerisch bedenklich ist.
\end{karte}

%%%%%%%%%%% Kapitel 7
\stupser{tief}{Bisektion halbiert das Intervall: $2^{-n}\le10^{-10}$ gibt $n\ge34$. Newton verdoppelt die gültigen Stellen: $1, 2, 4, 8, 16$.}{22pt}
\begin{karte}{tief}{5.22}{\typtypisch}{\meta{2}{15 min}{Band 5, Kap.\ 7}}
\kartentitel{Bisektion gegen Newton}
\worum{Konvergenzordnung konkret: linear heißt eine feste Zahl Stellen pro Schritt, quadratisch heißt Verdopplung. Und wann Newton auf linear zurückfällt.}
\methoden{Bisektion, Newton-Verfahren, Konvergenzordnung}
\begin{enumerate}[label=(\alph*),leftmargin=1.6em,itemsep=2pt]
\item Wie viele Bisektionsschritte braucht man auf $[1, 2]$ für $\sqrt2$ mit Genauigkeit $10^{-10}$? Wie viele Newton-Schritte (Aufgabe 5.24)?
\item Newton auf $f(x) = x^{2}$ (doppelte Nullstelle bei $0$) ab $x_0 = 1$: Berechnen Sie $x_1, \dots, x_4$. Welche Konvergenzordnung?
\item Newton auf $f(x) = \sqrt[3]x$ ab $x_0 = 1$: Berechnen Sie $x_1, x_2, x_3$. Was passiert?
\item Formulieren Sie die Voraussetzungen, unter denen Newton quadratisch konvergiert.
\end{enumerate}
\end{karte}

\stupser{akzent}{$f(x) = x^{3} - 2x + 2$, $x_0 = 0$: $x_1 = 0 - \frac2{-2} = 1$, $x_2 = 1 - \frac11 = 0$.}{22pt}
\begin{karte}{akzent}{5.23}{\typstolper}{\meta{2}{15 min}{Band 5, Kap.\ 7}}
\kartentitel{Newton im Kreis}
\worum{Newton konvergiert quadratisch --- wenn er konvergiert. Zwei Beispiele, in denen er es nicht tut: ein Zyklus und ein Weglaufen ins Unendliche.}
\methoden{Newton-Verfahren, Gegenbeispiel, Dämpfung}
\begin{enumerate}[label=(\alph*),leftmargin=1.6em,itemsep=2pt]
\item $f(x) = x^{3} - 2x + 2$, $x_0 = 0$. Berechnen Sie $x_1, x_2, x_3, x_4$. Skizzieren Sie, was geometrisch passiert.
\item Wo liegt die reelle Nullstelle? Wählen Sie einen Startwert, von dem aus Newton konvergiert, und rechnen Sie drei Schritte.
\item $f(x) = \arctan x$, $x_0 = 1{,}5$. Berechnen Sie $x_1, x_2, x_3$. Was passiert, und ab welchem $|x_0|$?
\item Welche praktischen Maßnahmen sichern Newton ab? (Stichworte: Dämpfung, Bisektion zum Einfangen, Abbruch bei wachsendem $|f|$.)
\end{enumerate}
\end{karte}

\stupser{gruen}{$x_{k+1} = \frac12\bigl(x_k + \frac2{x_k}\bigr)$ --- das Heron-Verfahren. Die Zahl der richtigen Stellen verdoppelt sich: $1{,}5$; $1{,}417$; $1{,}41422$; $1{,}414213562$.}{22pt}
\begin{karte}{gruen}{5.24}{\typwerk}{\meta{1}{10 min}{Band 5, Kap.\ 7, Aufgabe 7.1}}
\kartentitel{$\sqrt2$ auf zehn Stellen}
\worum{Newton auf $x^{2} - 2$ ist das Verfahren, das die Babylonier kannten. In fünf Schritten von $1$ auf $16$ Stellen.}
\methoden{Newton-Verfahren, Konvergenzordnung}
Berechnen Sie $\sqrt2$ mit dem Newton-Verfahren auf zehn Stellen genau, ausgehend von $x_0 = 1$. Zählen Sie die richtigen Stellen nach jedem Schritt.
\end{karte}

\stupser{gruen}{Fixpunkt: $|g'(x)| = |\sin x|\le0{,}72$ nahe der Lösung --- linear, Fehler mal $0{,}67$ pro Schritt. Newton: drei Schritte.}{22pt}
\begin{karte}{gruen}{5.25}{\typwerk}{\meta{2}{15 min}{Band 5, Kap.\ 7, Aufgabe 7.2}}
\kartentitel{$x = \cos x$: Fixpunkt oder Newton}
\worum{Dieselbe Gleichung, zwei Verfahren: Die Fixpunktiteration braucht $30$ Schritte, Newton $4$. Der Fixpunktsatz sagt vorher, dass beide funktionieren.}
\methoden{Fixpunktiteration, Kontraktion, Newton-Verfahren}
Lösen Sie $x = \cos x$ mit der Fixpunktiteration $x_{k+1} = \cos x_k$ (Kontraktionsbedingung prüfen!) und mit Newton, jeweils ab $x_0 = 1$. Vergleichen Sie die Schrittzahlen für $10^{-6}$.
\end{karte}

\stupser{gruen}{$F(x, y) = (x^{2} + y^{2} - 4,\ e^{x} + y - 1)$, Jacobi-Matrix $\begin{pmatrix}2x & 2y\\ e^{x} & 1\end{pmatrix}$. Ein Newton-Schritt ist ein lineares $2\times2$-System.}{22pt}
\begin{karte}{gruen}{5.26}{\typwerk}{\meta{3}{20 min}{Band 5, Kap.\ 7, Aufgabe 7.3}}
\kartentitel{Newton in zwei Dimensionen}
\worum{Statt durch $f'$ zu teilen, löst man mit der Jacobi-Matrix. Sonst ändert sich nichts --- auch nicht die quadratische Konvergenz.}
\methoden{Newton-Verfahren, Jacobi-Matrix}
Lösen Sie $x^{2} + y^{2} = 4$, $e^{x} + y = 1$ mit dem Newton-Verfahren ab $(1;\ -1{,}7)$. Zwei Schritte von Hand; wie viele Stellen stimmen danach?
\end{karte}

%%%%%%%%%%% Kapitel 8
\stupser{tief}{Newton-Schema: $f[1, 4] = \frac13$, $f[4, 9] = \frac15$, $f[1, 4, 9] = \frac{1/5 - 1/3}{8} = -\frac1{60}$. Fehlerformel $\frac{f'''(\xi)}{3!}\,\omega(x)$.}{22pt}
\begin{karte}{tief}{5.27}{\typtypisch}{\meta{2}{15 min}{Band 5, Kap.\ 8}}
\kartentitel{Der Interpolationsfehler}
\worum{Ein Interpolationspolynom aufstellen (Lagrange und Newton), auswerten, den Fehler messen und mit der Schranke vergleichen --- die oft viel zu pessimistisch ist.}
\methoden{Lagrange-Interpolation, Newton-Schema, Fehlerschranke}
$f(x) = \sqrt x$ an den Knoten $1, 4, 9$.
\begin{enumerate}[label=(\alph*),leftmargin=1.6em,itemsep=2pt]
\item Stellen Sie $p_2$ in Lagrange- und in Newton-Form auf.
\item Berechnen Sie $p_2(2)$ und den Fehler gegenüber $\sqrt2$.
\item Berechnen Sie die Fehlerschranke $\frac{\max|f'''|}{6}|\omega(2)|$ mit $\omega(x) = (x - 1)(x - 4)(x - 9)$. Warum ist sie so grob?
\item An welcher Stelle in $[1, 9]$ ist $|\omega|$ am größten? Was heißt das für die Wahl der Knoten?
\end{enumerate}
\end{karte}

\stupser{akzent}{$f(x) = \frac1{1 + 25x^{2}}$ auf $[-1, 1]$: Bei $11$ gleichabständigen Knoten ist der maximale Fehler $1{,}9$, bei $21$ Knoten $60$. Die Funktion hat Werte zwischen $0$ und $1$.}{22pt}
\begin{karte}{akzent}{5.28}{\typstolper}{\meta{2}{15 min}{Band 5, Kap.\ 8}}
\kartentitel{Höherer Grad, schlechteres Ergebnis}
\worum{Das Runge-Phänomen: Mehr Stützstellen machen das Interpolationspolynom nicht besser, sondern am Rand katastrophal --- wenn die Stützstellen gleichabständig sind.}
\methoden{Runge-Phänomen, Gegenbeispiel, Tschebyschow-Knoten}
$f(x) = \dfrac1{1 + 25x^{2}}$ auf $[-1, 1]$.
\begin{enumerate}[label=(\alph*),leftmargin=1.6em,itemsep=2pt]
\item Interpolieren Sie $f$ an $x = -1, 0, 1$ (Grad $2$) und an $x = -1, -\frac12, 0, \frac12, 1$ (Grad $4$). Werten Sie beide bei $x = 0{,}9$ aus.
\item Die maximalen Fehler für gleichabständige Knoten sind: Grad $4$: $0{,}44$; Grad $10$: $1{,}9$; Grad $20$: $60$. Für Tschebyschow-Knoten: $0{,}40$; $0{,}11$; $0{,}015$. Deuten Sie.
\item Woran liegt es? Betrachten Sie $\omega(x) = \prod(x - x_i)$ nahe dem Rand und die Ableitungen von $f$ (Pole bei $\pm\frac i5$).
\item Was tut man stattdessen? Nennen Sie drei Auswege.
\end{enumerate}
\end{karte}

\stupser{gruen}{Dividierte Differenzen: $1, 7, 19$; dann $3, 6$; dann $1$. Das Polynom ist $1 + x + 3x(x - 1) + x(x - 1)(x - 2)$.}{22pt}
\begin{karte}{gruen}{5.29}{\typwerk}{\meta{1}{10 min}{Band 5, Kap.\ 8, Aufgabe 8.1}}
\kartentitel{Das Newton-Schema}
\worum{Dividierte Differenzen in Dreiecksform, dann das Polynom ablesen --- und erkennen, dass es $x^{3} + 1$ ist.}
\methoden{Newton-Schema}
Interpolieren Sie $(0, 1)$, $(1, 2)$, $(2, 9)$, $(3, 28)$ mit dem Newton-Schema und multiplizieren Sie aus.
\end{karte}

\stupser{gruen}{$x_k = \cos\frac{(2k + 1)\pi}{10}$, $k = 0, \dots, 4$: $\pm0{,}951$, $\pm0{,}588$, $0$. Für diese Knoten ist $\max|\omega| = 2^{-4}$.}{22pt}
\begin{karte}{gruen}{5.30}{\typwerk}{\meta{2}{15 min}{Band 5, Kap.\ 8, Aufgabe 8.2}}
\kartentitel{Tschebyschow-Knoten}
\worum{Die Knoten, die $\max|\omega|$ minimieren: Nullstellen des Tschebyschow-Polynoms. Hier für fünf Punkte, mit Vergleich.}
\methoden{Tschebyschow-Knoten, Minimax-Eigenschaft}
Bestimmen Sie die Tschebyschow-Knoten für $n = 4$ (fünf Punkte) auf $[-1, 1]$ und vergleichen Sie $\max|\omega|$ mit dem für gleichabständige Punkte.
\end{karte}

%%%%%%%%%%% Kapitel 9
\stupser{gruen}{Momentengleichungen mit $h = 1$: $M_{i-1} + 4M_i + M_{i+1} = 6(y_{i+1} - 2y_i + y_{i-1})$, $M_0 = M_3 = 0$. Rechte Seiten $-18$ und $18$.}{22pt}
\begin{karte}{gruen}{5.31}{\typwerk}{\meta{2}{20 min}{Band 5, Kap.\ 9, Aufgabe 9.1}}
\kartentitel{Ein natürlicher Spline}
\worum{Vier Punkte, drei kubische Stücke, zwei unbekannte Momente. Der Spline ist das, was ein biegsames Lineal durch die Punkte macht.}
\methoden{Kubischer Spline, Tridiagonales System}
Bestimmen Sie den natürlichen kubischen Spline durch $(0, 1)$, $(1, 2)$, $(2, 0)$, $(3, 1)$. Geben Sie die drei Polynomstücke an und prüfen Sie die Glattheit bei $x = 1$.
\end{karte}

\stupser{gruen}{$N = 4$, $\omega = i$: $c_k = \frac14\sum_jy_ji^{-kj}$. Mit $i^{-1} = -i$, $i^{-2} = -1$, $i^{-3} = i$ ist das reine Buchhaltung.}{22pt}
\begin{karte}{gruen}{5.32}{\typwerk}{\meta{2}{10 min}{Band 5, Kap.\ 9, Aufgabe 9.2}}
\kartentitel{Eine DFT von Hand}
\worum{Vier Zahlen rein, vier komplexe Zahlen raus. Einmal selbst gerechnet, damit man weiß, was die FFT so schnell tut.}
\methoden{Diskrete Fourier-Transformation, Parseval}
Berechnen Sie die DFT der Folge $y = (1, 2, 1, 0)$ von Hand (mit Normierung $\frac1N$) und deuten Sie die Koeffizienten.
\end{karte}

%%%%%%%%%%% Kapitel 10
\stupser{gruen}{Legendre: $c_0 = \frac12\int_{-1}^{1}e^{x}\,dx$, $c_1 = \frac32\int_{-1}^{1}xe^{x}\,dx$. Das zweite Integral ist $\frac2e$.}{22pt}
\begin{karte}{gruen}{5.33}{\typwerk}{\meta{2}{15 min}{Band 5, Kap.\ 10, Aufgabe 10.1}}
\kartentitel{Beste Quadratmittel-Approximation}
\worum{Orthogonale Projektion auf die Polynome ersten Grades --- mit Legendre-Polynomen ohne Gleichungssystem. Und der Vergleich mit Taylor, der am Rand verliert.}
\methoden{Orthogonale Projektion, Legendre-Polynome}
Bestimmen Sie die beste Quadratmittel-Approximation von $f(x) = e^{x}$ durch ein Polynom ersten Grades auf $[-1, 1]$ und vergleichen Sie mit der Taylorentwicklung $1 + x$.
\end{karte}

\stupser{gruen}{Drei Alternanten: Ränder $0$, $1$ und der Scheitel von $r(x) = x^{2} - ax - b$. Aus $r(0) = r(1) = E$ folgt $a = 1$.}{22pt}
\begin{karte}{gruen}{5.34}{\typwerk}{\meta{2}{15 min}{Band 5, Kap.\ 10, Aufgabe 10.2}}
\kartentitel{Beste gleichmäßige Approximation}
\worum{Tschebyschows Alternantensatz an einem Beispiel, bei dem man alles ausrechnen kann: Der Fehler nimmt dreimal abwechselnd sein Maximum an.}
\methoden{Alternantensatz, Minimax-Eigenschaft}
Bestimmen Sie die beste gleichmäßige Approximation von $f(x) = x^{2}$ durch ein Polynom ersten Grades auf $[0, 1]$ und den Fehler $E$.
\end{karte}

%%%%%%%%%%% Kapitel 11
\stupser{tief}{Fehler bei $h$ und bei $h/2$ vergleichen: Verhältnis $2^{p}$, $p$ die Ordnung. Trapez $p = 2$, Simpson $p = 4$.}{22pt}
\begin{karte}{tief}{5.35}{\typtypisch}{\meta{2}{15 min}{Band 5, Kap.\ 11}}
\kartentitel{Die Fehlerordnung messen}
\worum{Die Ordnung eines Verfahrens kann man \emph{messen}: Schrittweite halbieren, Fehler vergleichen. Das ist der schnellste Test, ob eine Implementierung stimmt.}
\methoden{Konvergenzordnung, Romberg-Schema}
Für $I = \int_0^1e^{-x^{2}}\,dx = 0{,}7468241$ liefern Trapezregel und Simpson-Regel mit $m$ Teilintervallen:
\begin{center}\small
\begin{tabular}{lcccc}
$m$ & $2$ & $4$ & $8$ & $16$\\ \hline
Trapez & $0{,}7313703$ & $0{,}7429841$ & $0{,}7458656$ & $0{,}7465846$\\
Simpson & $0{,}7471804$ & $0{,}7468554$ & $0{,}7468261$ & $0{,}7468243$
\end{tabular}
\end{center}
\begin{enumerate}[label=(\alph*),leftmargin=1.6em,itemsep=2pt]
\item Berechnen Sie die Fehler und die Fehlerquotienten bei Halbierung. Welche Ordnungen ergeben sich?
\item Schätzen Sie aus zwei Trapezwerten den Fehler ohne Kenntnis von $I$ (Richardson).
\item Wie viele Teilintervalle brauchte die Trapezregel für $10^{-10}$? Simpson?
\item Warum ist der Simpson-Quotient erst bei feinen Gittern nahe $16$?
\end{enumerate}
\end{karte}

\stupser{gruen}{$h = \frac14$; Trapez $= h\bigl(\frac{f_0}2 + f_1 + f_2 + f_3 + \frac{f_4}2\bigr)$, Simpson $= \frac h3(f_0 + 4f_1 + 2f_2 + 4f_3 + f_4)$.}{22pt}
\begin{karte}{gruen}{5.36}{\typwerk}{\meta{2}{15 min}{Band 5, Kap.\ 11, Aufgabe 11.1}}
\kartentitel{Trapez und Simpson}
\worum{Die beiden Grundregeln mit fünf Funktionswerten --- und die Fehlerschätzung über die zweite bzw. vierte Ableitung.}
\methoden{Trapezregel, Simpson-Regel, Fehlerschranke}
Berechnen Sie $\int_0^1e^{-x^{2}}\,dx$ mit Trapez- und Simpson-Regel für $m = 4$ und schätzen Sie die Fehler mit den Fehlerformeln ab.
\end{karte}

\stupser{gruen}{$R_{k,j} = R_{k,j-1} + \frac{R_{k,j-1} - R_{k-1,j-1}}{4^{j} - 1}$. Die erste Spalte sind Trapezwerte, die zweite Simpson.}{22pt}
\begin{karte}{gruen}{5.37}{\typwerk}{\meta{2}{15 min}{Band 5, Kap.\ 11, Aufgabe 11.2}}
\kartentitel{Das Romberg-Schema}
\worum{Richardson-Extrapolation, systematisch: Aus vier Trapezwerten wird ein Ergebnis mit zehn Stellen.}
\methoden{Romberg-Schema}
Berechnen Sie $\int_0^1e^{-x^{2}}\,dx$ mit dem Romberg-Schema bis $R_{3,3}$ und geben Sie den Fehler jeder Spalte an.
\end{karte}

\stupser{gruen}{Knoten $t = 0, \pm\sqrt{3/5}$, Gewichte $\frac89, \frac59, \frac59$ auf $[-1, 1]$; transformieren mit $x = \frac{t + 1}2$, $dx = \frac{dt}2$.}{22pt}
\begin{karte}{gruen}{5.38}{\typwerk}{\meta{2}{15 min}{Band 5, Kap.\ 11, Aufgabe 11.3}}
\kartentitel{Gauß--Legendre mit drei Punkten}
\worum{Drei geschickt gewählte Punkte sind exakt für Polynome bis Grad $5$ --- Simpson mit drei Punkten nur bis Grad $3$. Der Unterschied im Beispiel: Faktor $40$.}
\methoden{Gauß-Quadratur}
Berechnen Sie $\int_0^1e^{-x^{2}}\,dx$ mit der Gauß--Legendre-Formel für drei Punkte und vergleichen Sie mit Simpson bei gleicher Punktzahl.
\end{karte}

%%%%%%%%%%% Kapitel 12
\stupser{akzent}{Expliziter Euler für $y' = -50y$: $y_{k+1} = (1 - 50h)y_k$. Mit $h = 0{,}1$ ist der Faktor $-4$.}{22pt}
\begin{karte}{akzent}{5.39}{\typstolper}{\meta{2}{15 min}{Band 5, Kap.\ 12}}
\kartentitel{Euler explodiert}
\worum{Eine Differentialgleichung, deren Lösung in Sekundenbruchteilen abklingt, und ein Verfahren, das stattdessen exponentiell wächst. Steifigkeit: der Grund, warum es implizite Verfahren gibt.}
\methoden{Steifigkeit, Stabilitätsgebiet, Impliziter Euler}
$y' = -50y$, $y(0) = 1$ (exakt $y = e^{-50t}$).
\begin{enumerate}[label=(\alph*),leftmargin=1.6em,itemsep=2pt]
\item Expliziter Euler mit $h = 0{,}1$: Berechnen Sie $y_1, \dots, y_4$ und vergleichen Sie mit der exakten Lösung.
\item Für welche $h$ bleibt der explizite Euler stabil? Wie viele Schritte braucht man dann für $[0, 1]$?
\item Impliziter Euler $y_{k+1} = y_k + hf(t_{k+1}, y_{k+1})$ mit $h = 0{,}1$: Berechnen Sie $y_1, \dots, y_4$.
\item Warum kann man beim impliziten Euler $h$ beliebig groß wählen? Was zahlt man dafür?
\end{enumerate}
\end{karte}

\stupser{gruen}{Euler: $y_1 = 0{,}5 + 0{,}2\cdot1{,}5 = 0{,}8$. RK4: vier Steigungen $k_1, \dots, k_4$, gewichtet $1, 2, 2, 1$.}{22pt}
\begin{karte}{gruen}{5.40}{\typwerk}{\meta{2}{20 min}{Band 5, Kap.\ 12, Aufgabe 12.1}}
\kartentitel{Euler und Runge--Kutta}
\worum{Zwei Schritte mit jedem Verfahren --- und ein Fehler, der beim einen in der ersten, beim anderen in der fünften Stelle liegt.}
\methoden{Euler-Verfahren, Runge–Kutta, Konvergenzordnung}
Lösen Sie $y' = y - t^{2} + 1$, $y(0) = 0{,}5$ auf $[0;\ 0{,}4]$ mit $h = 0{,}2$ mit Euler und mit RK4. Exakt: $y(t) = (t + 1)^{2} - \frac12e^{t}$.
\end{karte}

%%%%%%%%%%% Kapitel 13
\stupser{gruen}{$-u'' \approx\frac{-u_{i-1} + 2u_i - u_{i+1}}{h^{2}}$ mit $h = \frac14$: tridiagonales System $16\begin{pmatrix}2 & -1 & 0\\ -1 & 2 & -1\\ 0 & -1 & 2\end{pmatrix}u = f$.}{22pt}
\begin{karte}{gruen}{5.41}{\typwerk}{\meta{2}{15 min}{Band 5, Kap.\ 13, Aufgabe 13.1}}
\kartentitel{Ein Randwertproblem auf drei Punkten}
\worum{Finite Differenzen: Aus der Differentialgleichung wird ein lineares Gleichungssystem. Mit drei Punkten ist der Fehler fünf Prozent --- und man sieht, wie er mit $h^{2}$ fällt.}
\methoden{Finite Differenzen, Tridiagonales System, Symmetrie ausnutzen}
Lösen Sie $-u'' = \pi^{2}\sin(\pi x)$, $u(0) = u(1) = 0$ mit $n = 3$ inneren Gitterpunkten und vergleichen Sie mit $u(x) = \sin(\pi x)$.
\end{karte}

%%%%%%%%%%% Kapitel 14
\stupser{tief}{Für $f = x^{2} + 10y^{2}$ ist $\nabla f = (2x, 20y)$; exakte Schrittweite $\alpha = \frac{g^{\mathsf T}g}{g^{\mathsf T}Ag}$ mit $A = \operatorname{diag}(2, 20)$.}{22pt}
\begin{karte}{tief}{5.42}{\typtypisch}{\meta{2}{15 min}{Band 5, Kap.\ 14}}
\kartentitel{Gradientenverfahren im Tal}
\worum{Steilster Abstieg zickzackt durch lange Täler; die Konvergenzrate hängt von der Konditionszahl ab. Newton macht es in einem Schritt --- bei quadratischen Funktionen.}
\methoden{Gradientenverfahren, Liniensuche, Konditionszahl}
$f(x, y) = x^{2} + 10y^{2}$, Start $(10, 1)$.
\begin{enumerate}[label=(\alph*),leftmargin=1.6em,itemsep=2pt]
\item Führen Sie drei Schritte des Gradientenverfahrens mit exakter Liniensuche aus. Notieren Sie $f$ nach jedem Schritt.
\item Berechnen Sie die Quotienten aufeinanderfolgender $f$-Werte und vergleichen Sie mit $\bigl(\frac{\kappa - 1}{\kappa + 1}\bigr)^{2}$, $\kappa = 10$.
\item Führen Sie einen Newton-Schritt aus.
\item Was passiert mit dem Gradientenverfahren bei $f = x^{2} + 1000y^{2}$? Wie viele Schritte für Faktor $10^{-6}$?
\end{enumerate}
\end{karte}

\stupser{gruen}{Bei $(-1, 1)$ ist $y - x^{2} = 0$: Gradient $(-4, 0)$, Hesse-Matrix $\begin{pmatrix}802 & 400\\ 400 & 200\end{pmatrix}$. Schritt: $H s = -\nabla F$.}{22pt}
\begin{karte}{gruen}{5.43}{\typwerk}{\meta{2}{15 min}{Band 5, Kap.\ 14, Aufgabe 14.1}}
\kartentitel{Ein Newton-Schritt auf Rosenbrock}
\worum{Die Testfunktion der Optimierer: ein gekrümmtes Tal mit Minimum bei $(1, 1)$. Ein Newton-Schritt landet in $x$ richtig --- und in $y$ weit daneben.}
\methoden{Newton-Verfahren, Hesse-Matrix, Dämpfung}
Führen Sie für $F(x, y) = (1 - x)^{2} + 100(y - x^{2})^{2}$ einen Newton-Schritt ab $(-1, 1)$ aus. Vergleichen Sie $F$ vorher und nachher. Was lehrt das über ungedämpftes Newton?
\end{karte}

\stupser{gruen}{Trefferanteil $\hat p\approx\frac\pi4$ mit Standardabweichung $\sqrt{p(1 - p)/N}$; $\hat\pi = 4\hat p$ hat Standardabweichung $4\sqrt{p(1 - p)/N}$.}{22pt}
\begin{karte}{gruen}{5.44}{\typwerk}{\meta{2}{10 min}{Band 5, Kap.\ 14, Aufgabe 14.2}}
\kartentitel{$\pi$ mit Monte-Carlo}
\worum{Zufallspunkte im Quadrat, Anteil im Viertelkreis: Die Methode ist einfach, konvergiert wie $\frac1{\sqrt N}$ --- und das ist in hohen Dimensionen ihre Stärke, in einer ihre Schwäche.}
\methoden{Monte-Carlo, Standardfehler, Simpson-Regel}
Schätzen Sie $\pi$ mit Monte-Carlo und bestimmen Sie, wie viele Punkte für drei sichere Nachkommastellen nötig sind. Vergleichen Sie mit der Simpson-Regel für $4\int_0^1\sqrt{1 - x^{2}}\,dx$.
\end{karte}
